% Part: computability
% Chapter: recursive-functions
% Section: pr-functions-computable

\documentclass[../../../include/open-logic-section]{subfiles}

\begin{document}

\olfileid{cmp}{rec}{cmp}
\olsection{الدوال العودية البدائية قابلة للحساب}

افترض أن دالة $h$ معرفة بالعودية البدائية:
\begin{eqnarray*}
h(\vec x, 0)     & = & f(\vec x) \\
h(\vec x, y + 1) & = & g(\vec x, y, h(\vec x, y))
\end{eqnarray*}
وافترض أن الدالتين $f$ و$g$ قابلتان للحساب. (نستعمل $\vec x$ اختصارًا
لـ$x_0$، …،~$x_{k-1}$.) من الواضح عندئذ إمكان حساب $h(\vec x,0)$،
لأنها ليست إلا $f(\vec x)$ التي نفترض قابليتها للحساب. ويمكن بعد ذلك
حساب $h(\vec x,1)$ أيضًا، لأن $1=0+1$، ولذلك ليست $h(\vec x,1)$ إلا
\begin{align*}
 h(\vec x, 1) & = g(\vec x, 0, h(\vec x, 0)) =  g(\vec x, 0, f(\vec x)).
\intertext{ويمكننا المواصلة على هذا النحو وحساب}
h(\vec x, 2) & = g(\vec x, 1, h(\vec x, 1)) = g(\vec x, 1, g(\vec x, 0, f(\vec x)))\\
h(\vec x, 3) & = g(\vec x, 2, h(\vec x, 2)) = g(\vec x, 2, g(\vec x, 1, g(\vec x, 0, f(\vec x))))\\
h(\vec x, 4) & = g(\vec x, 3, h(\vec x, 3)) = g(\vec x, 3, g(\vec x, 2, g(\vec x, 1, g(\vec x, 0, f(\vec x)))))\\
& \vdots
\end{align*}
ومن ثم، لحساب $h(\vec x,y)$ بوجه عام، نحسب تباعًا $h(\vec x,0)$ و
$h(\vec x,1)$ و…، حتى نبلغ $h(\vec x,y)$.

وهكذا ينتج التعريف بالعودية البدائية دالة جديدة قابلة للحساب إذا كانت
الدالتان $f$ و$g$ قابلتين للحساب. وينتج تركيب الدوال أيضًا دالة قابلة
للحساب إذا كانت $f$ ودوال $g_i$ قابلة للحساب.

ولما كانت الدوال الأساسية $\Zero$ و$\Succ$ و$\Proj{n}{i}$ قابلة
للحساب، وكان التركيب والعودية البدائية ينتجان دوال قابلة للحساب من
دوال قابلة للحساب، فإن كل دالة عودية بدائية قابلة للحساب.

\end{document}
